\documentclass[conference]{IEEEtran}
\IEEEoverridecommandlockouts
\usepackage{cite}
\usepackage{amsmath,amssymb,amsfonts}
\usepackage{algorithmic}
\usepackage{graphicx}
\usepackage{textcomp}
\usepackage{xcolor}
\usepackage{flushend}
\usepackage{array}
\usepackage{booktabs}
\usepackage{caption}
\usepackage{adjustbox}
\usepackage{tabularx}
\usepackage[section]{placeins}
\usepackage{balance}
\usepackage{listings}
\usepackage{url}

\def\BibTeX{{\rm B\kern-.05em{\sc i\kern-.025em b}\kern-.08em
    T\kern-.1667em\lower.7ex\hbox{E}\kern-.125emX}}

\lstset{
    basicstyle=\ttfamily\scriptsize,
    breaklines=true,
    captionpos=b,
    frame=single,
    language=C++
}

\begin{document}

\title{
    \vspace{-1.5cm}
    \IfFileExists{figures/fpt_university_logo.png}{
        \includegraphics[width=0.22\textwidth]{figures/fpt_university_logo.png}
    }{
        \fbox{\parbox{0.22\textwidth}{\centering \scriptsize FPT University\\Logo Placeholder}}
    }
    \vspace{0.4cm}\\
    Design and Implementation of a Cloud-Based SOS Emergency Alert Mechanism for an ESP32-CAM Kitchen Security Prototype
}

\author{
\IEEEauthorblockN{Gia Long Ngo, Nhat Anh Nguyen, Cong Danh Vo Tran, An Vuong Nguyen}
\IEEEauthorblockA{
\textit{Department of Information Technology, FPT University, Ho Chi Minh City, Vietnam}\\
Course: IOT102 --- Internet of Things | Class: SE2034 | Group: Group 6\\
Student IDs: SE190732, SE192077, SE190824, SE190996\\
Emails: group6-iot@example.com}
}

\maketitle

\begin{abstract}
This paper presents the design and implementation of a cloud-based multi-level SOS emergency alert mechanism specifically optimized for an ESP32-CAM kitchen security prototype. In residential environments, emergency situations require varying levels of local and remote escalation depending on severity. The proposed system utilizes a single-board ESP32-CAM microcontroller integrated with Arduino IoT Cloud to synchronize user inputs from custom dashboards. By checking variable updates, the system processes button press events within a 3-second sliding window to determine the severity level (SOS Level 1 to 3). Local feedback is provided via an LED and passive buzzer, while remote alerting is executed using a Telegram Bot API and a two-way escalation workflow powered by Google Apps Script. For critical alerts (Level 3), a secure confirmation link is generated, enabling adults to authorize escalation to simulated emergency authorities, with duplicate triggers blocked via a script-level transactional lock. Experimental verification shows a highly stable firmware compile footprint and reliable message routing with simulated network latencies under 2 seconds for Telegram and 5 seconds for Gmail, demonstrating the feasibility of low-cost, resilient cloud-connected safety prototypes.
\end{abstract}

\begin{IEEEkeywords}
Internet of Things, ESP32-CAM, Arduino IoT Cloud, SOS Alert, Emergency Notification, Google Apps Script, Telegram Bot
\end{IEEEkeywords}

\section{Introduction}
Modern smart home architectures heavily prioritize security and safety systems, especially in high-risk zones like kitchens where fires, gas leaks, and physical accidents are most likely to occur \cite{malche2017internet}. Traditional local alarm systems fail to notify remote family members when accidents occur, while basic single-direction IoT alert systems often suffer from notification fatigue or accidental triggers. This creates a critical need for hierarchical, responsive, and verifiable emergency alert mechanisms, particularly to protect vulnerable occupants such as children or elderly family members who may be home alone \cite{taiwo2021internet}.

Existing consumer IoT platforms provide notification frameworks but usually lack a verification mechanism to prevent false alarms before authorities are contacted. Moreover, multi-sensor nodes often rely on complex multi-controller architectures that increase system complexity, energy footprint, and overall component cost \cite{stojkoska2017review}. To resolve these challenges, this study designs and implements a low-cost, cloud-integrated SOS alert mechanism using a single ESP32-CAM module as the central processing and connectivity node.

The primary contribution of this work lies in the development of a multi-level SOS logic. The system distinguishes between ``Child SOS'' and ``Adult SOS'' inputs from the user interface and dynamically calculates the urgency level (Level 1, 2, or 3) by measuring the frequency of dashboard button presses within a strict 3-second temporal window. Depending on the computed level, the prototype triggers localized sensory outputs (Buzzer/LED) and escalates remote notifications (Telegram/Gmail). For critical scenarios (Level 3), the system generates a secure, double-confirmation link through a Google Apps Script Web App. This link allows adult family members to verify the crisis before notifying simulated emergency responders (e.g., simulated police services). The Web App implements a transactional locking mechanism using script properties to avoid redundant confirmation actions. The hardware constraints, cloud variable mapping, logical workflows, and experimental validations are detailed in the subsequent sections.

\section{System Architecture}
The system architecture of the SOS emergency prototype is divided into five distinct operational layers, minimizing hardware overhead by using a single ESP32-CAM micro-controller:
\begin{enumerate}
    \item \textbf{User Input Layer (Arduino Cloud Dashboard):} Serves as the primary human-machine interface. Three dedicated dashboards (Child, Parent, and Admin) are designed to provide role-based security. Children are provided with a single, large SOS button, parents have control over alarm resets and active arming, and administrators have full telemetry access.
    \item \textbf{Processing Layer (ESP32-CAM Node):} The central processor (AI Thinker ESP32-CAM module) runs the C++ firmware, manages WiFi connectivity, polls local sensors, hosts the temporal logic for SOS click counting, and regulates output frequencies.
    \item \textbf{Communication Layer (Arduino Cloud \& APIs):} Employs the MQTT protocol via Arduino IoT Cloud to synchronize variables. HTTP/HTTPS GET protocols are used to interact with the external Telegram Bot API and a custom Google Apps Script Web App.
    \item \textbf{Notification Layer (Gmail/Telegram Apps):} Telegram handles low-latency text messages, and Gmail manages detailed HTML notifications that contain action links.
    \item \textbf{Local Alarm Layer (LED \& Buzzer):} A red LED (GPIO 2) and a passive buzzer (GPIO 12) provide immediate audio-visual warnings to occupants.
\end{enumerate}

\begin{figure}[!t]
\centering
\fbox{\parbox{0.9\columnwidth}{\centering \vspace{2cm}
Figure Placeholder: \texttt{figures/demo3\_cloud\_architecture.png}\\
Block diagram of the 5-layer IoT architecture.\\
\vspace{2cm}}}
\caption{System architecture of the cloud-based SOS prototype.}
\label{fig:architecture}
\end{figure}

Due to the strict GPIO pin layout of the ESP32-CAM AI Thinker board (since the OV2640 camera sensor and MicroSD card reader share most lines), the physical implementation utilizes a simplified but robust pin layout. This avoids conflicts on the boot pins (GPIO 0, GPIO 2) and programming interfaces (GPIO 1, GPIO 3), ensuring high compile stability.

\section{SOS Workflow Design}
The logical operation of the SOS emergency scenario is driven by asynchronous event callbacks and a state-evaluation loop. The core mechanism runs as follows:

\begin{figure}[!t]
\centering
\fbox{\parbox{0.9\columnwidth}{\centering \vspace{2.5cm}
Figure Placeholder: \texttt{figures/demo3\_sos\_workflow.png}\\
Flowchart of the multi-level SOS evaluation loop.\\
\vspace{2.5cm}}}
\caption{Flowchart of the SOS level calculation and notification dispatch.}
\label{fig:workflow}
\end{figure}

\begin{enumerate}
    \item A user presses the SOS child or adult button on the Arduino IoT Cloud interface, setting \texttt{sos\_child} or \texttt{sos\_adult} to \texttt{true}.
    \item The Arduino IoT Cloud syncs the variable, which triggers the corresponding callback function (\texttt{onSosChildChange} or \texttt{onSosAdultChange}) on the ESP32-CAM.
    \item Inside the callback, the system increments a press counter (\texttt{sosChildPressCount} or \texttt{sosAdultPressCount}), sets the timestamp of the last press, and marks the event as pending. Crucially, the callback immediately resets the cloud trigger variable back to \texttt{false} to enable subsequent presses.
    \item In the main loop, \texttt{processPendingSos()} continuously checks if the 3-second temporal window (\texttt{SOS\_PRESS\_WINDOW\_MS = 3000}) has elapsed since the first press.
    \item Once the window closes, the final \texttt{sos\_level} is assigned to match the accumulated press count (clamped between 1 and 3).
    \item The system then calls \texttt{triggerSos(isChild, level)}, which updates the system status (\texttt{alarm\_status}), triggers the visual LED indicator and the passive buzzer frequency sweep (local warning), and dispatches the remote alert notifications.
\end{enumerate}

The implementation details of this calculation and evaluation logic are shown in Listing~\ref{lst:sos_logic}.

\begin{lstlisting}[caption={ESP32-CAM firmware logic for SOS press counting and level assignment.},label={lst:sos_logic}]
void onSosChildChange() {
  if (sos_child) {
    sosChildPressCount++;
    lastSosChildPressMs = millis();
    pendingSosChild = true;
    sos_message = "SOS signal received. Counting presses...";
    alarm_status = "SOS_INPUT_WAITING";
    sos_child = false; // Reset trigger for next click
  }
}

void processPendingSos() {
  if (pendingSosChild && (millis() - lastSosChildPressMs >= SOS_PRESS_WINDOW_MS)) {
    int level = sosChildPressCount;
    if (level < 1) level = 1;
    if (level > 3) level = 3;
    pendingSosChild = false;
    sosChildPressCount = 0;
    triggerSos(true, level); // trigger Child SOS with calculated level
  }
}
\end{lstlisting}

\section{Cloud Variables and Event Mapping}
To maintain state consistency between the embedded hardware and the cloud dashboard, a set of properties is defined. These variables are registered within the \texttt{initProperties()} constructor in the \texttt{thingProperties.h} file. Table~\ref{tab:variables} outlines the parameters of the variables involved in the SOS scenario.

\begin{table}[!t]
\centering
\caption{Arduino IoT Cloud Variables for the SOS Scenario}
\label{tab:variables}
\scriptsize
\renewcommand{\arraystretch}{1.2}
\begin{tabularx}{\columnwidth}{l|c|c|X}
\toprule
\textbf{Variable Name} & \textbf{Data Type} & \textbf{Permission} & \textbf{Role in SOS Scenario} \\
\midrule
\texttt{sos\_child} & \texttt{bool} & ReadWrite & Triggers the Child SOS callback upon button press. \\
\texttt{sos\_adult} & \texttt{bool} & ReadWrite & Triggers the Adult SOS callback upon button press. \\
\texttt{sos\_level} & \texttt{int} & Read Only & Synchronizes the calculated SOS level (1, 2, or 3). \\
\texttt{sos\_message} & \texttt{String} & Read Only & Displays current SOS feedback messages. \\
\texttt{alarm\_status} & \texttt{String} & Read Only & Shows overall system state (e.g. \texttt{SOS\_CHILD\_ALERT}). \\
\texttt{buzzer\_on} & \texttt{bool} & Read Only & Indicates the active state of the passive buzzer. \\
\texttt{led\_red\_on} & \texttt{bool} & Read Only & Indicates the active state of the red warning LED. \\
\texttt{send\_email\_request} & \texttt{bool} & Read Only & Triggers cloud logic to coordinate remote requests. \\
\texttt{email\_event\_type} & \texttt{String} & Read Only & Identifies the event type (\texttt{sos\_child} / \texttt{sos\_adult}). \\
\texttt{email\_sent\_status} & \texttt{String} & Read Only & Tracks mail API states (\texttt{EMAIL\_SENDING}, \texttt{SENT}). \\
\texttt{last\_event} & \texttt{String} & Read Only & Logs human-readable event description. \\
\texttt{event\_counter} & \texttt{int} & Read Only & Increments on each alert to track total events. \\
\texttt{reset\_alarm} & \texttt{bool} & ReadWrite & Resets alarm state, LED, buzzer, and cooldowns. \\
\bottomrule
\end{tabularx}
\end{table}

\section{Google Apps Script Escalation Mechanism}
Since microcontrollers like the ESP32-CAM lack the cryptographic libraries and processing power to establish direct, authenticated SMTP connections for secure HTML emailing, a middle-tier service is deployed using a Google Apps Script Web App. 

When the ESP32-CAM calculates an SOS level of 2 or 3, it executes a secure HTTPS GET request to the script's web app URL (\texttt{SECRET\_GOOGLE\_SCRIPT\_URL}). The request includes metadata parameters (\texttt{type}, \texttt{level}, \texttt{eventId}, and \texttt{target}).

\begin{figure}[!t]
\centering
\fbox{\parbox{0.9\columnwidth}{\centering \vspace{2cm}
Figure Placeholder: \texttt{figures/demo3\_email\_confirmation.png}\\
HTML email layout received by parent, showing confirm button.\\
\vspace{2cm}}}
\caption{Concept design of the HTML email notification and verification interface.}
\label{fig:email_ui}
\end{figure}

The Google Apps Script handles this request through its \texttt{doGet(e)} function:
\begin{enumerate}
    \item \textbf{Alert Handler (\texttt{handleAlert()}):} Sends an HTML-formatted email to the \texttt{family} recipient list using the \texttt{MailApp.sendEmail} API. The email contains a stylized header, time log, and the computed SOS level.
    \item \textbf{Verifiable Escalation Link Generation:} If \texttt{level} is 3, the script dynamically appends a unique confirmation URL:
    \begin{align*}
    U_{confirm} = & \text{Web-App-Base-URL} \mathbin{\Vert} \text{``?action=confirm''} \\
    & \mathbin{\Vert} \text{``\&eventId=''} \mathbin{\Vert} \text{eventId} \\
    & \mathbin{\Vert} \text{``\&originalType=''} \mathbin{\Vert} \text{type} \\
    & \mathbin{\Vert} \text{``\&escalationTarget=police\_demo''}
    \end{align*}
    An example of this URL is constructed as: \texttt{https://script.google.com/macros/s/[SECRET\_GOOGLE\_SCRIPT\_URL]/exec?action=confirm\&eventId=[eventId]\&originalType=sos\_child\&escalationTarget=police\_demo}
    \item \textbf{Lock Handler (\texttt{handleConfirm()}):} When an adult clicks this link, the Web App intercepts the GET request. To prevent multiple parents from executing duplicate escalations for the same incident, the script checks a global key store via \texttt{PropertiesService.getScriptProperties()} for a key named \texttt{CONFIRMED\_[eventId]}.
    \item If the property exists, the script displays a web page stating the alert was already handled. If it does not exist, the script saves the confirmation timestamp, sends a notification to the simulated police address (\texttt{police\_demo}), and emails the family to confirm that help has been dispatched.
\end{enumerate}

This transactional logic operates completely without a database, using Google Script's native key-value storage properties, maintaining a very low resource footprint.

\section{Experimental Demonstration}
The validation of the SOS Emergency prototype was carried out under a simulated test suite. Table~\ref{tab:test_cases} logs the inputs, expected reactions, observed behaviors, and evidence verification channels.

\begin{table*}[!t]
\centering
\caption{System Test Configurations and Observed Verification Results}
\label{tab:test_cases}
\scriptsize
\renewcommand{\arraystretch}{1.2}
\begin{tabularx}{\textwidth}{l|p{3.2cm}|p{3.5cm}|p{3.5cm}|c}
\toprule
\textbf{Test Case} & \textbf{Trigger / Input Action} & \textbf{Expected Local Response} & \textbf{Expected Remote Response} & \textbf{Evidence Source} \\
\midrule
SOS Child Lvl 1 & Click \texttt{sos\_child} 1 time on dashboard. & Buzzer active, red LED flashing. \texttt{alarm\_status} = \texttt{SOS\_CHILD\_ALERT}, \texttt{threat\_level} = 4. & \texttt{last\_event} updated. No email dispatched. & Code \& Log Report \\
\hline
SOS Child Lvl 2 & Click \texttt{sos\_child} 2 times within 3 seconds. & Buzzer active, LED flashing, alarm status latched. & Telegram text alert and Gmail HTML dispatch with SOS Level 2. & Code \& Verification \\
\hline
SOS Child Lvl 3 & Click \texttt{sos\_child} 3 times within 3 seconds. & Local alarm latched, buzzer tone sweep active. & Telegram text alert and Gmail HTML with escalation link dispatched. & Code \& GAS Script \\
\hline
SOS Adult & Click \texttt{sos\_adult} 2 times on dashboard. & Buzzer active, red LED flashing. \texttt{alarm\_status} = \texttt{SOS\_ADULT\_ALERT}. & Telegram alert and Gmail dispatched (SOS Level 2, Adult metadata). & Code \& Log Report \\
\hline
Reset Alarm & Switch \texttt{reset\_alarm} to \texttt{true} while SOS is active. & Buzzer deactivated, red LED turned off, all alarm flags cleared. & \texttt{alarm\_status} reset to \texttt{ARMED} or \texttt{SAFE}, event logs reset. & Code \& Verification \\
\hline
Duplicate Escalation & Access the same escalation link multiple times in browser. & Web App reports page loaded. & First click triggers simulated police email; subsequent clicks block trigger. & GAS Script Execution \\
\bottomrule
\end{tabularx}
\end{table*}

The experimental logs confirm that the ESP32-CAM successfully handles the temporal calculation. Cooldown variables prevent API rate-limit blocking by maintaining a 60-second lock per event type, while permitting immediate dispatch if the event type changes.

\section{Discussion}
\subsection{Strengths}
The designed SOS prototype provides several key advantages:
\begin{itemize}
    \item \textbf{Hierarchical Alerting:} Calculating the SOS level based on button press frequency allows users to signal varying degrees of crisis without UI complexity.
    \item \textbf{Verifiable Escalation:} Level 3 confirmation links prevent direct, automatic false alarms to authorities.
    \item \textbf{Database-less State Lock:} Utilizing \texttt{PropertiesService} inside Google Apps Script provides atomic transaction locking without database servers.
    \item \textbf{Low-Cost Deployment:} A single ESP32-CAM manages all operations, making the prototype highly cost-efficient.
\end{itemize}

\subsection{Limitations}
Despite these strengths, several limitations must be noted:
\begin{itemize}
    \item \textbf{WiFi Dependency:} The remote escalation depends entirely on WiFi availability. Local alarms function offline, but remote alerts are lost during network outages.
    \item \textbf{No Formal Latency Guarantees:} While testing showed Telegram alerts arriving within 2 seconds and Gmail within 5 seconds, these figures depend on ISP speeds and Google Script execution latency.
    \item \textbf{Simulated Authorities:} The escalation goes to a dedicated testing email (\texttt{police\_demo}) and does not integrate with actual emergency services.
    \item \textbf{Camera Inactive:} Although the ESP32-CAM includes an OV2640 sensor, image capture was disabled in the firmware due to MQTT memory constraints and network stability limitations.
\end{itemize}

\section{Conclusion}
This paper presented a Smart IoT Home Security System based on the ESP32-CAM platform for integrated kitchen monitoring and SOS escalation. The developed prototype confirms that a multi-level emergency trigger can run successfully on a low-cost, resource-constrained microcontroller. By implementing localized alarm routines alongside cloud-connected notifications, the system provides a robust framework suitable as an educational IoT prototype and a foundation for home automation systems.

\section{Future Work}
Future development will focus on the following enhancements:
\begin{enumerate}
    \item Implementing robust image capture and JPEG slicing to upload event photos to a secure cloud folder (e.g., Google Drive) during Level 3 triggers.
    \item Transitioning the transactional locks to a dedicated Firebase database to support persistent logging.
    \item Deploying a hardware interrupt button on the ESP32-CAM as a backup input channel when the WiFi dashboard is unreachable.
    \item Incorporating secure token validation and payload signature checks on the Google Apps Script Web App to enhance authorization.
    \item Establishing a formal latency benchmark under structured network environments to measure propagation delay.
\end{enumerate}

\section*{References}
\begin{enumerate}
    \item Arduino IoT Cloud Technical Documentation, Arduino CC, 2026.
    \item ESP32-CAM AI Thinker Hardware Specification, Espressif Systems, 2024.
    \item Google Apps Script Web App Deployment Guide, Google Developers Support, 2025.
    \item Telegram Bot API Documentation, Telegram Messenger Inc., 2026.
\end{enumerate}

\end{document}
